特别地, 可以利用这一定理得到以下结果.

\begin{Theorem}
\label{thm:hilbert-schmidt-kernel-compactness}
设核 $k:\Omega\times\Omega\to\mathbb{R}$ 为 Hilbert-Schmidt 的. 那么由下式定义的积分算子 $K\in\mathcal{L}(L^2(\Omega))$ 为紧算子:\[K(u)(\bm{x}):=\int_\Omega k(\bm{x},\bm{y})u(\bm{y})\,\mathrm{d}\bm{y}\]
\end{Theorem}

\begingroup
\postdisplaypenalty=10000
\renewcommand{\endProof}{\par\nobreak\hfill$\blacksquare$\par\vspace{0.1cm}}
\begin{Proof}
设 $\{\phi_i(\bm{x})\}$ 为 $L^2(\Omega)$ 的一组标准正交基. 那么利用 Section~\ref{subsec:distribution-direct-product} 中的方法, 可以证明 $\{\phi_i(\bm{x})\phi_j(\bm{y})\}$ 为 $L^2(\Omega\times\Omega)$ 的一组基. 将 $k$ 关于这组基展开, 得到
\begin{equation}\label{eq:hilbert-schmidt-kernel-product-basis-expansion}
k(\bm{x},\bm{y})=\sum_{i,j=1}^\infty k_{ij}\phi_i(\bm{x})\phi_j(\bm{y})
\end{equation}
其中级数在 $L^2(\Omega\times\Omega)$ 范数下收敛, 且
\begin{equation}\label{eq:hilbert-schmidt-kernel-product-basis-coefficients}
k_{ij}:=\int_\Omega\int_\Omega k(\bm{x},\bm{y})\phi_i(\bm{x})\phi_j(\bm{y})\,\mathrm{d}\bm{x}\,\mathrm{d}\bm{y}.
\end{equation}
此外, 由 \eqref{eq:hilbert-basis-parseval-characterization} 有
\begin{equation}\label{eq:hilbert-schmidt-kernel-coefficient-square-sum}
\int_\Omega\int_\Omega|k(\bm{x},\bm{y})|^2\,\mathrm{d}\bm{x}\,\mathrm{d}\bm{y}=\sum_{i,j=1}^\infty|k_{ij}|^2.
\end{equation}
现在用下式定义算子 $K_n\in\mathcal{L}(L^2(\Omega))$:
\begin{equation}\label{eq:hilbert-schmidt-truncated-kernel-operator}
K_n(u)(\bm{x}):=\int_\Omega k_n(\bm{x},\bm{y})u(\bm{y})\,\mathrm{d}\bm{y}
\end{equation}
其中
\begin{equation}\label{eq:hilbert-schmidt-rectangular-kernel-truncation}
k_n(\bm{x},\bm{y})=\sum_{i,j=1}^n k_{ij}\phi_i(\bm{x})\phi_j(\bm{y}).
\end{equation}
分别将 $k_n$ 和 $K_n$ 称为\textbf{可分离核}和\textbf{可分离算子}. 容易看出, 可分离算子具有有限秩, 因而是紧算子.

现在利用 Lemma~\ref{lem:hilbert-schmidt-kernel-boundedness} 中的方法, 得到
\begin{equation}\label{eq:hilbert-schmidt-truncation-operator-norm-bound}
\|K-K_n\|^2\leq\int_\Omega\int_\Omega|k(\bm{x},\bm{y})-k_n(\bm{x},\bm{y})|^2\,\mathrm{d}\bm{x}\,\mathrm{d}\bm{y}.
\end{equation}
再利用 \eqref{eq:hilbert-schmidt-kernel-coefficient-square-sum}, 得到
\begin{equation}\label{eq:hilbert-schmidt-truncation-coefficient-tail}
\begin{aligned}
&\lim_{n\to\infty}\int_\Omega\int_\Omega|k(\bm{x},\bm{y})-k_n(\bm{x},\bm{y})|^2\,\mathrm{d}\bm{x}\,\mathrm{d}\bm{y}\\
&\qquad=\lim_{n\to\infty}\sum_{i,j=n+1}^\infty|k_{ij}|^2=0.
\end{aligned}
\end{equation}
因而 $K_n$ 在算子范数下收敛到 $K$, 所以 Theorem~\ref{thm:norm-limit-of-compact-operators} 表明 $K$ 为紧算子.
\end{Proof}
\endgroup
